\section{A complete norm from reciprocal-integer samples (NS-78)}
\label{sec:ns78-sampling}

\paragraph{Purpose.}
NS-75 identifies the signed divisor defects but shows that cancelling
individual jumps need not decrease the norm. Here we retain the entire
residual through its reciprocal-integer values, with constants independent
of the number and size of the coefficients. This is an exact geometric
reduction and a complete physical-space replay, not the missing arithmetic
estimate. The infinite sample sum is not replaced by a finite prefix.

For arbitrary real $\alpha,c_1,\ldots,c_N$, put
\[
 R=\alpha\tau-\sum_{n\leq N}c_n\sigma_n,\qquad
 \eta=\sum_{n\leq N}\frac{c_n}{n},\qquad
 z_m=R(1/m),\qquad w_m=\frac1{m(m+1)}.
\]
Define the complete nonnegative arithmetic quantity
\begin{equation}
 \mathcal S(c)=\eta^2+
       \sum_{m=1}^\infty w_m(z_m^2+z_{m+1}^2).
 \label{eq:ns78-sampling-functional}
\end{equation}
Its samples are explicitly
\[
 z_m=-\eta m+\sum_{k\leq m}
 \left(\alpha\mathbf1_{k=1}+\sum_{\substack{n\mid k\\n\leq N}}c_n\right)
                    \log(m/k).
\]
All cancellation inside each displayed sample is retained before squaring.

\begin{proposition}[Uniform complete sampling comparison]
\label{prop:ns78-sampling}
For every real $\alpha$ and finite real coefficient vector,
\begin{equation}
 \frac1{16}\mathcal S(c)\leq\|R\|_2^2
                 \leq\frac{21}{16}\mathcal S(c).
 \label{eq:ns78-sampling-comparison}
\end{equation}
Both sides include the exterior tail and every reciprocal cell. The
infinite series in~\eqref{eq:ns78-sampling-functional} converges.
\end{proposition}
\begin{proof}
Change variables to $t=1/x$ and write $r(t)=R(1/t)$, with measure
$dt/t^2$. On $0<t<1$, $r(t)=-\eta t$, so the complete exterior
energy is exactly $\eta^2$. By NS-75, on $m\leq t\leq m+1$,
\[
 r(t)=-\eta t+D_m\log t-L_m,
\quad D_m=\sum_{k\leq m}\delta_k,
\quad L_m=\sum_{k\leq m}\delta_k\log k.
\]
Here $\delta_k=\alpha\mathbf1_{k=1}+\sum_{n\mid k,\,n\leq N}c_n$
is the signed divisor defect.
Let $h_m=\log(1+1/m)$, $\theta=\log(t/m)/h_m$, and set
\[
 p(t)=(1-\theta)z_m+\theta z_{m+1},\qquad
 b(t)=\theta-(t-m).
\]
Exactly $r(t)=p(t)+\eta b(t)$ on every cell.
Under $t=m e^{h_m\theta}$ the measure is
$(h_m/m)e^{-h_m\theta}d\theta$. Its density divided by $w_m$
is between $m h_m$ and $(m+1)h_m$, hence between $1/2$ and $2$.
Also
\[
 \int_0^1((1-\theta)a+\theta d)^2d\theta
       =\frac{a^2+ad+d^2}{3}
       \in\left[\frac{a^2+d^2}{6},\frac{a^2+d^2}{2}\right].
\]
Thus, writing $V=\int_1^\infty p(t)^2dt/t^2$ and
$S=\sum_{m\geq1}w_m(z_m^2+z_{m+1}^2)$, we get
$S/12\leq V\leq S$, initially for finite unions of cells.

Concavity of $\log(m+v)$ for $0\leq v\leq1$, and its second
derivative bound, give
\[
 0\leq b(m+v)\leq\frac{v(1-v)}{2m^2h_m}
       \leq\frac1{8m^2h_m}\leq\frac1{4m}.
\]
Consequently
\[
 \Gamma:=\int_1^\infty b(t)^2\frac{dt}{t^2}
 \leq\frac1{16}\sum_{m\geq1}\frac{w_m}{m^2}
 \leq\frac1{16},
\]
since $\sum w_m=1$. The original finite combination lies in $L^2$,
so $p=r-\eta b$ is square integrable. Monotone passage through the
finite-cell estimates now proves finiteness of $S$ and their complete
versions; no convergence of $S$ was assumed.

Cauchy--Schwarz and Young's inequality give
\[
 |2\eta\langle p,b\rangle|
       \leq\eta^2/4+4\Gamma V\leq(\eta^2+V)/4.
\]
Using the exact identity
$\|R\|_2^2=\eta^2+V+2\eta\langle p,b\rangle+\eta^2\Gamma$,
we obtain
\[
 \frac34(\eta^2+V)\leq\|R\|_2^2
 \leq\frac{21}{16}\eta^2+\frac54V
 \leq\frac{21}{16}(\eta^2+V).
\]
Inserting $S/12\leq V\leq S$ proves the proposition.
\end{proof}

\paragraph{The arithmetic input in explicit coordinates.}
Specialize to $\alpha=1$. By the fixed-order smoothed Nyman--Beurling
criterion in NS-61, RH is
equivalent to the existence of coefficient vectors, of unbounded allowed
length, with $\mathcal S(c)\to0$. For the actual optimized coefficients,
this is equivalent to $E_N\to0$. Boundedness of $\mathcal S(c)$ alone
does not suffice; it is not the bounded-Weil-floor criterion.
The local geometric comparison is unconditional, but no sequence with
vanishing complete $\mathcal S$ has been produced. A finite collection
of small samples cannot replace its infinite sum or its cofinal limit.

The case $\alpha=0$ also bounds the complete energy of any finite
correction vector. Its comparison quantity still contains all signed
arithmetic samples; this is not an all-vector comparison with the
smooth repaired model excluded in NS-70.

\paragraph{Exact inversion and the remaining arithmetic constraint.}
The samples also recover the coefficients without a Gram solve. Put
$D_0=0$. Continuity at every reciprocal knot gives
\[
 \eta=-z_1,\qquad
 D_m=\frac{z_{m+1}-z_m+\eta}{\log(1+1/m)},\qquad
 \delta_m=D_m-D_{m-1}.
\]
M\"obius inversion then gives, extending $c_n=0$ beyond $N$,
\begin{equation}
 c_n=\sum_{d\mid n}\mu(n/d)\delta_d-\alpha\mu(n).
 \label{eq:ns78-sample-inverse}
\end{equation}
For $n>N$ the right-hand side must therefore vanish. These infinitely
many arithmetic compatibility conditions distinguish admissible sample
sequences from arbitrary small knot values.

In particular, norm convergence of residuals with $\alpha=1$ forces
$c_n\to-\mu(n)$ at each fixed index $n$. Indeed
$|z_m|\leq4\sqrt{m(m+1)}\,\|R\|_2$ follows from the complete
sampling comparison, while $|\eta|\leq\|R\|_2$ follows from the
exterior tail. The finite formulas above then force $D_m,\delta_m\to0$
for every fixed $m$, and~\eqref{eq:ns78-sample-inverse} gives the claim.
This necessary coefficientwise limit supplies neither a uniform growing-index
bound nor a norm-convergent sharp M\"obius truncation; the latter was
excluded in NS-62/66. No interchange of those two limits is made.

\subsection{Complete cell integrals and explicit tails}
\label{subsec:ns78-tails}
For $m\geq1$, put $a=\log m$, $d=\log(m+1)$ and
\[
 I_{1,m}=\frac{a+1}{m}-\frac{d+1}{m+1},\qquad
 I_{2,m}=\frac{a^2+2a+2}{m}-\frac{d^2+2d+2}{m+1}.
\]
The exact cell contribution is
\begin{equation}
 \int_m^{m+1}\frac{r(t)^2}{t^2}dt
 =\eta^2+D_m^2I_{2,m}+L_m^2w_m
 -\eta D_m(d^2-a^2)+2\eta L_m(d-a)-2D_mL_mI_{1,m}.
 \label{eq:ns78-exact-cell}
\end{equation}
The signed cross terms are evaluated before accumulating cell energies.

We also need a complete tail, since $m$ extends indefinitely even when
$c$ has finite support. Write
\[
 u=\alpha-\tfrac12\sum c_n,\qquad
 v=\tfrac12\sum c_n\log n-\tfrac12\log(2\pi)\sum c_n,
 \qquad D=\tfrac16\sum n|c_n|.
\]
For $t\geq N$,
\begin{equation}
 r(t)=u\log t+v+e(t),\qquad |e(t)|\leq D/t.
 \label{eq:ns78-complete-asymptotic}
\end{equation}
Indeed $A'(z)=\{z\}/z$ almost everywhere and Stirling's formula at
integer endpoints identifies the limiting constant
$A(z)-\tfrac12\log z\to\tfrac12\log(2\pi)$.
Using the periodic Bernoulli polynomials,
\[
 A(z)-\tfrac12\log(2\pi z)
 =-\int_z^\infty\frac{B_1(\{t\})}{t}dt
 =\frac{B_2(\{z\})}{2z}
       -\frac12\int_z^\infty\frac{B_2(\{t\})}{t^2}dt.
\]
The first integral is an improper convergent integral; the second is
absolutely convergent. Since $|B_2|\leq1/6$, the remainder is at most
$1/(6z)$. The periodic $B_2$ is continuous at integer joins, so no
endpoint mass is discarded. Summing this bound gives
\eqref{eq:ns78-complete-asymptotic}.

Fix an integer $M\geq N$ and put $q=u\log M+v$. The complete
main-tail energy and remainder budget are
\[
 T_M=\frac{q^2+2uq+2u^2}{M},\qquad
 U_M=\frac{D^2}{3M^3}.
\]
Thus the actual integral from $M$ to infinity differs from $T_M$ by
at most $2\sqrt{T_MU_M}+U_M$. For the infinite sample tail, set
$Q=|q|+(D+|u|)/M$. Then
\begin{equation}
 \sum_{m\geq M}w_m(z_m^2+z_{m+1}^2)
 \leq \frac{2(Q^2+2|u|Q+2u^2)}{M}.
 \label{eq:ns78-sampling-tail}
\end{equation}
To verify this last bound, on each $m\leq t\leq m+1$ both endpoint
values have absolute value at most $Q+|u|\log(t/M)$:
use~\eqref{eq:ns78-complete-asymptotic} and
$\log((m+1)/t)\leq1/M$. Integrating the square with measure $dt/t^2$
and summing all cells proves~\eqref{eq:ns78-sampling-tail}.
These are bounds on the complete tails, not estimates based on a zero
prefix or an unmeasured remainder.

\paragraph{Independent complete physical-space check.}
For the actual optimized coefficients at $N=16,64,256$, the interval
calculation accumulates every cell through $M=65536$ and encloses the
entire remaining physical and sample tails by the formulas above.
At both 256 and 384 bits, the resulting complete physical norm encloses
the independently evaluated full Gram norm. The complete ratios satisfy
\[
 \mathcal S(c_{16})/E_{16}\in(3.323,3.325),\qquad
 \mathcal S(c_{64})/E_{64}\in(2.9002,2.9004),\qquad
 \mathcal S(c_{256})/E_{256}\in(2.644,2.646).
\]
The analytic physical-tail uncertainty divided by $E_N$ is below
$1.64\cdot10^{-6}$, $5.20\cdot10^{-6}$ and $2.38\cdot10^{-4}$,
respectively. Replaying $N=256$ with $M=131072$ and the doubled complete
Gram-kernel cutoff reduces that tail ratio below $4.17\cdot10^{-5}$.
All solves retain Arb balls. The precise finite part, main tail and
remainder radius are archived separately; the printed combined decimal
balls may be wider due to decimal serialization. Thirty exact Maxima
checks cover the cell, mass, interpolation and tail identities.
This is an independent norm evaluation on finite coefficient vectors,
not a cofinal estimate of their sample sums.
